\chapter{La desintegración radiactiva}\label{ch:g12:radioactive-decay}

En el sótano de un museo, un contador Geiger crepita sobre una astilla de madera del
sarcófago de un faraón. Cada chasquido carece de ley: nada lo anunció, nada predice el
siguiente. Y sin embargo, para el viernes el laboratorio imprimirá una fecha con un
siglo de precisión. Este capítulo resuelve la paradoja --- cómo unos sucesos
perfectamente aleatorios uno a uno suman, por billones, los relojes más regulares que
tenemos.

\section{Sin ley uno a uno, relojería por moles}

El curso pasado (\cref{ch:g11:nucleus-radioactivity}) conocimos los \omterm{def:g11:fundamental-interactions:ladder}{núcleos}
inestables, ajustamos sus ecuaciones $\alpha$ y $\beta$ y vimos desaparecer la mitad
de cualquier muestra grande en cada \omterm{def:g11:nucleus-radioactivity:halflife}{periodo de semidesintegración} $T_{1/2}$. Lo que no
sabíamos decir era \emph{cuándo entre escalón y escalón} --- ni por qué un proceso
aleatorio respeta horario alguno. Las dos preguntas tienen una sola respuesta.

\begin{definition}[Constante de desintegración]\label{def:g12:radioactive-decay:lambda}
Para cada \omterm{def:g11:nucleus-radioactivity:notation}{nucleido} inestable existe una constante $\lambda$, su \emph{constante de
desintegración}\index{constante de desintegración} (\omterm{def:g10:orders-of-magnitude:unit}{unidad} \unit{s^{-1}}), tal que
durante cualquier intervalo corto $\dd t$ cada \omterm{def:g11:fundamental-interactions:ladder}{núcleo} tiene la probabilidad
$p = \lambda\,\dd t$ de desintegrarse --- la misma para todos los \omterm{def:g11:fundamental-interactions:ladder}{núcleos} de la
especie, sea cual sea su edad, e insensible a la temperatura, la \omterm{def:g10:pressure:pressure}{presión} o la química.
\end{definition}

\begin{remark}[Sin ley por separado, exactos en multitud]\label{rem:g12:radioactive-decay:crowds}
Un \omterm{def:g11:fundamental-interactions:ladder}{núcleo} que lleva esperando diez mil años tiene exactamente la misma probabilidad de
desintegrarse este segundo que otro fabricado esta mañana: no se desgasta, no avisa
--- cuál será el siguiente es genuinamente impredecible. Con las multitudes ocurre
otra cosa. Lanza $100$ monedas: espera $50$ caras, pero no te sorprendas ante $43$ ---
fluctuaciones del orden de $\sqrt n$, aquí un $10\,\%$. Lanza $10^{20}$: la
fluctuación relativa $1/\sqrt n$ vale $10^{-10}$, y el resultado es seguro hasta la
décima cifra. Un gramo de materia alberga unos $10^{22}$ \omterm{def:g11:fundamental-interactions:ladder}{núcleos}: la falta de ley se
promedia hasta convertirse en relojería, una ley exacta para $N(t)$, el número de
\omterm{def:g11:fundamental-interactions:ladder}{núcleos} todavía intactos en el instante $t$.
\end{remark}

\begin{omfigure}
\begin{tikzpicture}
\begin{axis}[omaxis, width=8.8cm, height=5.0cm,
    xlabel={$t$ (en $T_{1/2}$)}, ylabel={supervivientes},
    xmin=0, xmax=4.3, ymin=0, ymax=21, xtick={0,1,2,3,4},
    legend pos=north east, legend style={font=\small}]
  \addplot[omDef, thick, domain=0:4.2, samples=60] {20*exp(-0.6931*x)};
  \addplot[omProp, only marks, mark=x] coordinates
    {(0,20) (0.5,17) (1,11) (1.5,8) (2,6) (2.5,4) (3,3) (3.5,2) (4,1)};
  \addplot[omExo, only marks, mark=o] coordinates
    {(0,20) (0.5,14) (1,9) (1.5,7) (2,4) (2.5,4) (3,2) (3.5,1) (4,1)};
  \legend{esperado, muestra 1, muestra 2}
\end{axis}
\end{tikzpicture}

{\small Dos muestras reales de $20$ \omterm{def:g11:fundamental-interactions:ladder}{núcleos} bajan a trompicones e impredeciblemente en
torno a la curva esperada; un mol de \omterm{def:g11:fundamental-interactions:ladder}{núcleos} se le pegaría hasta la décima cifra.}
\end{omfigure}

\section{La ley de desintegración}

\begin{proposition}[La ecuación de evolución]\label{prop:g12:radioactive-decay:ode}
En una muestra grande de $N$ \omterm{def:g11:fundamental-interactions:ladder}{núcleos}, durante $\dd t$ el número varía en
\[
\dd N = -\lambda N\,\dd t, \qquad\text{es decir,}\qquad \frac{\dd N}{\dd t} = -\lambda N :
\]
la población decrece a un ritmo proporcional a sí misma.
\end{proposition}

\begin{proof}
Cada uno de los $N$ \omterm{def:g11:fundamental-interactions:ladder}{núcleos} se desintegra durante $\dd t$ con probabilidad
$\lambda\,\dd t$: desintegraciones esperadas, $N\lambda\,\dd t$, y para $N$ grande el
recuento real se pega al esperado (\cref{rem:g12:radioactive-decay:crowds}); cada
desintegración retira un \omterm{def:g11:fundamental-interactions:ladder}{núcleo}.
\end{proof}

\begin{theorem}[Ley de desintegración radiactiva]\label{thm:g12:radioactive-decay:law}
Una muestra que alberga $N_0$ \omterm{def:g11:fundamental-interactions:ladder}{núcleos} en $t = 0$ alberga, en el instante $t$,
\[
N(t) = N_0 \, e^{-\lambda t}.
\]
\end{theorem}

\begin{proof}
Comprobación por sustitución:
$\dd N/\dd t = N_0(-\lambda)\,e^{-\lambda t} = -\lambda N(t)$, y
$N(0) = N_0 e^0 = N_0$: la ecuación y el arranque, ambos correctos. (Que ninguna
\emph{otra} función encaja se demuestra, con integrales, en el volumen
universitario.)
\end{proof}

\begin{definition}[Decrecimiento exponencial]\label{def:g12:radioactive-decay:exponential}
Una magnitud que sigue $N_0\,e^{-\lambda t}$ experimenta un \emph{decrecimiento
exponencial}\index{decrecimiento exponencial}: en tiempos iguales se pierden
\emph{fracciones} iguales. Es exactamente el \omterm{def:g12:rc-rl-circuits:capacitor}{condensador} que se descarga del
\cref{ch:g12:rc-rl-circuits}, con $\lambda$ haciendo el papel de $1/(RC)$: una sola
ecuación diferencial, aprendida una vez, gobierna por igual circuitos y \omterm{def:g11:fundamental-interactions:ladder}{núcleos}.
\end{definition}

\begin{proposition}[Periodo de semidesintegración]\label{prop:g12:radioactive-decay:halflife}
El \omterm{def:g11:nucleus-radioactivity:halflife}{periodo de semidesintegración}
(\cref{def:g11:nucleus-radioactivity:halflife}) de un \omterm{def:g11:nucleus-radioactivity:notation}{nucleido} lo fija su \omterm{def:g12:radioactive-decay:lambda}{constante de
desintegración}: $T_{1/2} = \ln 2/\lambda$, y tras $n$ \omterm{def:g12:oscillators-and-time:oscillator}{periodos}
$N = N_0/2^{\,n}$, sea $n$ entero o no.
\end{proposition}

\begin{proof}
$N(T_{1/2}) = N_0/2$ exige $e^{-\lambda T_{1/2}} = \tfrac12$, es decir,
$\lambda T_{1/2} = \ln 2$; y $N(n\,T_{1/2}) = N_0\,e^{-n\ln 2} = N_0/2^{\,n}$ para
cualquier $n$: la exponencial enhebra la escalera y rellena lo que hay entre escalón y
escalón.
\end{proof}

\begin{omfigure}
\begin{tikzpicture}
\begin{axis}[omaxis, width=8.8cm, height=5.2cm,
    xlabel={$t$ (en $T_{1/2}$)}, ylabel={$N/N_0$},
    xmin=0, xmax=5, ymin=0, ymax=1.08, xtick={0,1,2,3,4,5},
    ytick={0.125,0.25,0.5,1}, yticklabels={$\tfrac18$,$\tfrac14$,$\tfrac12$,$1$}]
  \addplot[omDef, thick, domain=0:5, samples=80] {exp(-0.6931*x)};
  \addplot[omExo, dashed] coordinates {(0,0.5) (1,0.5) (1,0)};
  \addplot[omExo, dashed] coordinates {(0,0.25) (2,0.25) (2,0)};
  \addplot[omExo, dashed] coordinates {(0,0.125) (3,0.125) (3,0)};
\end{axis}
\end{tikzpicture}

{\small $N(t) = N_0\,e^{-\lambda t}$ pasa por todos los escalones de la escalera --- y
ahora dice además qué ocurre entre ellos.}
\end{omfigure}

\begin{example}[Veinte potencias de diez]\label{ex:g12:radioactive-decay:range}
Con $1$ año $= \qty{3.156e7}{s}$:
\begin{center}\small
\begin{tabular}{llll}
\omterm{def:g11:nucleus-radioactivity:notation}{nucleido} & $T_{1/2}$ & $\lambda$ (\unit{s^{-1}}) & \\
polonio-214 & \qty{1.6e-4}{s} & \num{4.2e3} & eslabón de la cadena del uranio\\
tecnecio-99m & $6.0$ horas & \num{3.2e-5} & diagnóstico por imagen\\
yodo-131 & $8.0$ días & \num{1.0e-6} & medicina del tiroides\\
cesio-137 & $30$ años & \num{7.3e-10} & lluvia radiactiva de reactor\\
carbono-14 & \num{5730} años & \num{3.8e-12} & arqueología\\
potasio-40 & \num{1.25e9} años & \num{1.8e-17} & datación de rocas\\
uranio-238 & \num{4.5e9} años & \num{4.9e-18} & calor interno de la Tierra\\
\end{tabular}
\end{center}
Una sola ley, veinte potencias de diez de $\lambda$.
\end{example}

\section{La actividad: contar los chasquidos}

\begin{proposition}[Actividad]\label{prop:g12:radioactive-decay:activity}
La \omterm{def:g11:nucleus-radioactivity:activity}{actividad} de una muestra --- sus desintegraciones por segundo, en \omterm{def:g11:nucleus-radioactivity:activity}{becquerelios}
(\cref{def:g11:nucleus-radioactivity:activity}) --- vale
\[
\mathcal A = \lambda N, \qquad\text{luego}\qquad \mathcal A(t) = \mathcal A_0 \, e^{-\lambda t}:
\]
proporcional a la población superviviente e igualmente exponencial.
\end{proposition}

\begin{proof}
$N$ \omterm{def:g11:fundamental-interactions:ladder}{núcleos}, cada uno con probabilidad $\lambda$ por segundo: $\lambda N$
desintegraciones por segundo; multiplicar el
\cref{thm:g12:radioactive-decay:law} por $\lambda$ da
$\mathcal A(t) = \mathcal A_0 e^{-\lambda t}$.
\end{proof}

\begin{example}[$N$ enorme, $\lambda$ minúscula]\label{ex:g12:radioactive-decay:body}
Un plátano crepita a unos \qty{15}{Bq}, un cuerpo humano a unos \qty{8000}{Bq} y un
\omterm{def:g10:orders-of-magnitude:unit}{metro} cúbico de granito a unos \qty{e6}{Bq} (\cref{ch:g11:nucleus-radioactivity}).
Detrás de esos modestos chasquidos hay poblaciones astronómicas: el potasio-40 del
cuerpo aporta $\mathcal A \approx \qty{4.4e3}{Bq}$ con
$\lambda = \qty{1.8e-17}{s^{-1}}$: $N = \mathcal A/\lambda \approx \num{2.5e20}$
\omterm{def:g11:fundamental-interactions:ladder}{núcleos}, de los cuales apenas unos pocos miles disparan cada segundo.
\end{example}

\begin{method}[Medir un periodo de semidesintegración]\label{met:g12:radioactive-decay:semilog}
\begin{enumerate}
  \item Cuenta desintegraciones con un contador Geiger a lo largo de intervalos
        iguales sucesivos: eso muestrea $\mathcal A$ en instantes sucesivos.
  \item Representa $\ln \mathcal A$ frente a $t$: como
        $\ln \mathcal A = \ln \mathcal A_0 - \lambda t$, un \omterm{def:g12:radioactive-decay:exponential}{decrecimiento exponencial}
        aparece como una \emph{recta} --- unos datos torcidos significan otra ley, o
        una mezcla de \omterm{def:g11:nucleus-radioactivity:notation}{nucleidos}.
  \item Lee $\lambda$ como menos la pendiente; $T_{1/2} = \ln 2/\lambda$.
\end{enumerate}
\end{method}

\begin{omfigure}
\begin{tikzpicture}
\begin{axis}[omaxis, width=8.8cm, height=5.0cm,
    xlabel={$t$ (\unit{s})}, ylabel={$\ln \mathcal A$},
    xmin=0, xmax=260, ymin=0, ymax=6.8, xtick={0,60,120,180,240}, ytick={2,4,6}]
  \addplot[omDef, thick, domain=0:250, samples=2] {5.99-0.00958*x};
  \addplot[omProp, only marks, mark=*, mark size=1.6pt] coordinates
    {(0,5.99) (60,5.42) (120,4.84) (180,4.26) (240,3.69)};
  \draw[omExo, dashed] (axis cs:60,5.42) -| (axis cs:180,4.26);
  \node[font=\small, above right] at (axis cs:110,4.28) {pendiente $= -\lambda$};
\end{axis}
\end{tikzpicture}

{\small Recuentos Geiger de un \omterm{def:g11:nucleus-radioactivity:notation}{nucleido} de vida corta: en ejes semilogarítmicos la
exponencial se convierte en una recta; aquí la pendiente da
$\lambda \approx \qty{9.6e-3}{s^{-1}}$, luego $T_{1/2} \approx \qty{72}{s}$.}
\end{omfigure}

\begin{remark}[Los becquerelios no son daño]\label{rem:g12:radioactive-decay:dose}
El \omterm{def:g11:nucleus-radioactivity:activity}{becquerelio} cuenta desintegraciones, no daños. El daño se sigue con la \omterm{rem:g11:nucleus-radioactivity:dose}{dosis} en
\omterm{rem:g11:nucleus-radioactivity:dose}{sieverts} (\cref{ch:g11:nucleus-radioactivity}), que pondera la \omterm{def:g10:energy-conservation:energy}{energía} depositada por
\omterm{def:g10:orders-of-magnitude:unit}{kilogramo} de tejido con el poder destructivo de cada radiación: una gammagrafía de
tecnecio de \qty{5e8}{Bq} es rutinaria, y muchos menos \omterm{def:g11:nucleus-radioactivity:activity}{becquerelios} de un emisor
$\alpha$ inhalado no lo son; convertir \unit{Bq} en \unit{Sv} es materia
universitaria.
\end{remark}

\section{Datar: leer el tiempo en un cociente}

La ley de desintegración recorrida al revés es un reloj: mide la fracción
superviviente y la exponencial nombra el tiempo transcurrido.

\begin{method}[Datación por radiocarbono]\label{met:g12:radioactive-decay:carbon}
\begin{enumerate}
  \item \emph{En vida}: el tejido vivo intercambia carbono con la atmósfera y lleva su
        proporción de carbono-14 --- un \omterm{def:g11:fundamental-interactions:ladder}{átomo} de cada $10^{12}$, lo que da
        $\mathcal A_0 = 13.6$ desintegraciones por minuto y por gramo de carbono.
  \item \emph{La muerte pone en marcha el reloj}: cesa la incorporación y la
        proporción decrece con $T_{1/2} = \num{5730}$ años.
  \item Mide la \omterm{def:g11:nucleus-radioactivity:activity}{actividad} $\mathcal A$ de la muestra por gramo de carbono;
        $\mathcal A = \mathcal A_0 e^{-\lambda t}$ se invierte en
        $t = \frac{1}{\lambda}\ln(\mathcal A_0/\mathcal A)
        = T_{1/2}\,\ln(\mathcal A_0/\mathcal A)/\ln 2$.
\end{enumerate}
\end{method}

\begin{example}[Un hogar en una cueva]\label{ex:g12:radioactive-decay:hearth}
El carbón vegetal de un hogar enterrado da $5.0$ desintegraciones por minuto y por
gramo: $\mathcal A_0/\mathcal A = 13.6/5.0 = 2.72$, luego
$t = \num{5730} \times \ln(2.72)/\ln 2 \approx \num{8.3e3}$ años. El fuego se apagó
hace ocho mil años --- la propia madera es el archivo.
\end{example}

\begin{omfigure}
\begin{tikzpicture}
\begin{axis}[omaxis, width=8.8cm, height=5.0cm,
    xlabel={edad (miles de años)}, ylabel={$\mathcal A/\mathcal A_0$},
    xmin=0, xmax=32, ymin=0, ymax=1.08, xtick={0,10,20,30},
    ytick={0.25,0.37,0.5,1}, yticklabels={$\tfrac14$,$0.37$,$\tfrac12$,$1$}]
  \addplot[omDef, thick, domain=0:31, samples=80] {exp(-0.6931*x/5.73)};
  \addplot[omExo, dashed] coordinates {(0,0.37) (8.3,0.37) (8.3,0)};
  \fill[omThm] (axis cs:8.3,0.37) circle (1.6pt);
\end{axis}
\end{tikzpicture}

{\small La curva de datación: se entra por el cociente medido y se baja hasta la edad.
El punto es el hogar del \cref{ex:g12:radioactive-decay:hearth}: cociente $0.37$, edad
\num{8.3e3} años.}
\end{omfigure}

\begin{example}[Relojes para rocas]\label{ex:g12:radioactive-decay:rocks}
El carbono-14 enmudece más allá de unos \num{50000} años; la geología guarda relojes
más lentos. \emph{Potasio--argón}: el potasio-40 ($T_{1/2} = \num{1.25e9}$ años) se
desintegra en argón, un gas que escapa de la lava fundida pero queda atrapado en
cuanto la roca cristaliza --- la solidificación pone el reloj a cero.
\emph{Uranio--plomo}: el uranio-238 encabeza una cadena de desintegraciones que
termina en el plomo-206 estable (\cref{ch:g11:nucleus-radioactivity}), y el cociente
plomo/uranio en un cristal lo data. Los dos relojes coinciden en los minerales más
antiguos ($\approx \num{4.4e9}$ años) y, a través de los meteoritos, en la edad del
sistema solar.
\end{example}

\begin{remark}[Cadenas, y el gas del sótano]\label{rem:g12:radioactive-decay:radon}
Entre el uranio y el plomo, la cadena pasa por el radio y después por el radón-222, un
\emph{gas} radiactivo ($T_{1/2} = 3.8$ días). Nacido en el granito y en el suelo, se
filtra en los sótanos y se acumula; sus hijos, emisores $\alpha$, se desintegran en
los pulmones --- para la mayoría de la gente, la mayor \omterm{rem:g11:nucleus-radioactivity:dose}{dosis} natural, y la más barata
de reducir: abre la ventana del sótano.
\end{remark}

\section{Ejercicios}

\begin{exercise}[$\star$]\label{exo:g12:radioactive-decay:1}
El yodo-131 tiene $\lambda = \qty{1.0e-6}{s^{-1}}$. Calcula su \omterm{def:g11:nucleus-radioactivity:halflife}{periodo de
semidesintegración} en segundos y después en días.
\end{exercise}

\begin{exercise}[$\star$]\label{exo:g12:radioactive-decay:2}
El radón-222 tiene $T_{1/2} = 3.8$ días. Calcula $\lambda$ en \unit{s^{-1}} y después
la \omterm{def:g11:nucleus-radioactivity:activity}{actividad} de una muestra de \num{1.0e10} \omterm{def:g11:fundamental-interactions:ladder}{núcleos}.
\end{exercise}

\begin{exercise}[$\star$]\label{exo:g12:radioactive-decay:3}
Verifica por sustitución que $N(t) = N_0\,e^{-\lambda t}$ satisface
$\dd N/\dd t = -\lambda N$ y $N(0) = N_0$; calcula después $N(T_{1/2})/N_0$.
\end{exercise}

\begin{exercise}[$\star$]\label{exo:g12:radioactive-decay:4}
Una fuente de cesio-137 ($T_{1/2} = 30$ años) tiene hoy una \omterm{def:g11:nucleus-radioactivity:activity}{actividad} de
\qty{4.0e5}{Bq}. Calcula su \omterm{def:g11:nucleus-radioactivity:activity}{actividad} dentro de $90$ años y después dentro de $300$
años.
\end{exercise}

\begin{exercise}[$\star$]\label{exo:g12:radioactive-decay:5}
La representación semilogarítmica de $\ln \mathcal A$ frente a $t$ para una muestra de
tecnecio es una recta que pasa por $(0, 20.00)$ y $(\qty{10}{horas}, 18.85)$. Halla
$\lambda$ (en \unit{h^{-1}} y en \unit{s^{-1}}) y $T_{1/2}$.
\end{exercise}

\begin{exercise}[$\star\star$]\label{exo:g12:radioactive-decay:6}
Un gramo de radio-226 ($T_{1/2} = 1600$ años; masa molar \qty{226}{g/mol};
$N_A = \qty{6.02e23}{mol^{-1}}$): calcula el número de \omterm{def:g11:fundamental-interactions:ladder}{núcleos}, después $\lambda$ y
después la \omterm{def:g11:nucleus-radioactivity:activity}{actividad} --- la primera \omterm{def:g10:orders-of-magnitude:unit}{unidad} de \omterm{def:g11:nucleus-radioactivity:activity}{actividad} de la historia, el curio.
\end{exercise}

\begin{exercise}[$\star\star$]\label{exo:g12:radioactive-decay:7}
Dos estudiantes sellan muestras de $20$ y de \num{2.0e20} \omterm{def:g11:fundamental-interactions:ladder}{núcleos} de un mismo
\omterm{def:g11:nucleus-radioactivity:notation}{nucleido}. ¿Qué espera cada uno al cabo de un \omterm{def:g11:nucleus-radioactivity:halflife}{periodo de semidesintegración}? ¿De quién
es fiable la predicción? (Estima la fluctuación relativa $1/\sqrt n$ en cada caso.)
\end{exercise}

\begin{exercise}[$\star\star$]\label{exo:g12:radioactive-decay:8}
Un cuerpo alberga unos \qty{140}{g} de potasio, de los cuales el $0.012\,\%$ es
potasio-40 ($T_{1/2} = \num{1.25e9}$ años, masa molar \qty{40}{g/mol}). Calcula el
número de \omterm{def:g11:fundamental-interactions:ladder}{núcleos} de potasio-40 y después su \omterm{def:g11:nucleus-radioactivity:activity}{actividad}; compárala con los
\qty{8000}{Bq} totales del cuerpo.
\end{exercise}

\begin{exercise}[$\star\star$]\label{exo:g12:radioactive-decay:9}
Una paciente recibe \qty{500}{MBq} de tecnecio-99m ($T_{1/2} = 6.0$ horas). Calcula la
\omterm{def:g11:nucleus-radioactivity:activity}{actividad} $15$ horas después y después el tiempo que tarda en bajar de \qty{1}{MBq}.
\end{exercise}

\begin{exercise}[$\star\star$]\label{exo:g12:radioactive-decay:10}
El carbón vegetal de una fosa da $3.4$ desintegraciones por minuto y por gramo de
carbono; un hueso de otra da el $30\,\%$ del ritmo de lo vivo
($\mathcal A_0 = 13.6$ por minuto y por gramo, $T_{1/2} = \num{5730}$ años). Data
ambas muestras.
\end{exercise}

\begin{exercise}[$\star\star$]\label{exo:g12:radioactive-decay:11}
Se sella un sótano con una tanda de radón-222 dentro ($T_{1/2} = 3.8$ días). ¿Al cabo
de cuánto tiempo ha caído al $1\,\%$ la \omterm{def:g11:nucleus-radioactivity:activity}{actividad} de esa tanda? Los sótanos reales
siguen siendo radiactivos durante décadas: ¿de dónde sale el radón nuevo, y por qué la
ventilación debe ser permanente y no puntual?
\end{exercise}

\begin{exercise}[$\star\star\star$]\label{exo:g12:radioactive-decay:12}
Demuestra, a partir de $N = N_0 e^{-\lambda t}$, que
$t = T_{1/2}\ln(N_0/N)/\ln 2$. Un hueso conserva el $5.0\,\%$ de su carbono-14 de
cuando estaba vivo: dátalo.
\end{exercise}

\begin{exercise}[$\star\star\star$]\label{exo:g12:radioactive-decay:13}
En un cristal que no atrapó nada de argón al solidificarse, cada \omterm{def:g11:fundamental-interactions:ladder}{núcleo} de potasio-40
desintegrado deja un \omterm{def:g11:fundamental-interactions:ladder}{átomo} de argón en su sitio. Demuestra que
$N_{\mathrm{Ar}}/N_{\mathrm K} = e^{\lambda t} - 1$; una roca presenta un cociente de
$3.0$: halla su edad ($T_{1/2} = \num{1.25e9}$ años).
\end{exercise}

\begin{exercise}[$\star\star\star$]\label{exo:g12:radioactive-decay:14}
Un contador Geiger da, en $t = 0$, $60$, $120$, $180$ y $240$ segundos, ritmos de
$400$, $225$, $127$, $71$ y $40$ cuentas por segundo. Calcula $\ln \mathcal A$ en cada
instante, comprueba que están alineados y deduce $\lambda$ y $T_{1/2}$.
\end{exercise}

\begin{exercise}[$\star\star\star$]\label{exo:g12:radioactive-decay:15}
Para un solo \omterm{def:g11:fundamental-interactions:ladder}{núcleo}, la probabilidad de sobrevivir $n$ \omterm{def:g11:nucleus-radioactivity:halflife}{periodos de semidesintegración}
es $1/2^{\,n}$. Calcula la probabilidad de que (a) un \omterm{def:g11:fundamental-interactions:ladder}{núcleo} dado sobreviva $10$
\omterm{def:g12:oscillators-and-time:oscillator}{periodos}; (b) los $8$ \omterm{def:g11:fundamental-interactions:ladder}{núcleos} vigilados sobrevivan todos a un \omterm{def:g12:oscillators-and-time:oscillator}{periodo}. (c) De
\num{1000} \omterm{def:g11:fundamental-interactions:ladder}{núcleos}, ¿cuántos se esperan tras $3$ \omterm{def:g12:oscillators-and-time:oscillator}{periodos} --- y puede la ley decir
\emph{cuáles}? Concluye sobre qué es lo que predice la ley.
\end{exercise}

\section{Problema: La edad de las cosas}

\begin{problem}[{Problema de fin de semana --- la edad de las cosas: un laboratorio de
museo calibra el reloj del carbono, data una estatua y un papiro, aprende por qué
ningún reloj marca más allá de su esfera y empuja una roca volcánica mil millones de
años atrás}]
\label{pb:g12:radioactive-decay:1}
Una semana en el laboratorio de datación de un museo. Datos: el tejido vivo presenta
$\mathcal A_0 = 13.6$ desintegraciones de carbono-14 por minuto y por gramo de
carbono; $T_{1/2} = \num{5730}$ años (carbono-14) y \num{1.25e9} años (potasio-40);
$1$ año $= \qty{3.156e7}{s}$; $N_A = \qty{6.02e23}{mol^{-1}}$; masa molar del carbono
\qty{12}{g/mol}; los contadores resuelven alrededor de $0.2$ desintegraciones por
minuto y por gramo.

\medskip
\noindent\textbf{Parte I --- Lunes: calibrar el reloj.}
\begin{enumerate}
  \item ¿Por qué es constante la proporción de carbono-14 en el tejido vivo, y qué
        cambia con la muerte?
  \item Calcula $\lambda$ para el carbono-14 en \unit{s^{-1}}.
  \item Verifica por sustitución que $N_0 e^{-\lambda t}$ resuelve
        $\dd N/\dd t = -\lambda N$.
  \item A partir de $\mathcal A_0$, calcula el número de \omterm{def:g11:fundamental-interactions:ladder}{átomos} de carbono-14 por
        gramo de carbono en el tejido vivo.
  \item ¿Cuántos \omterm{def:g11:fundamental-interactions:ladder}{átomos} de carbono alberga un gramo? Deduce la proporción de
        carbono-14 --- compárala con ``uno de cada $10^{12}$''.
  \item Demuestra que una \omterm{def:g11:nucleus-radioactivity:activity}{actividad} medida $\mathcal A$ data una muestra en
        $t = T_{1/2}\,\ln(\mathcal A_0/\mathcal A)/\ln 2$.
\end{enumerate}

\medskip
\noindent\textbf{Parte II --- Martes: la estatua y el papiro.}
\begin{enumerate}[resume]
  \item Una estatua de madera da $9.1$ desintegraciones por minuto y por gramo. Data
        la madera.
  \item ¿Qué acontecimiento exactamente fecha eso --- la talla u otra cosa? ¿Qué
        precaución se sigue de ello?
  \item Un papiro egipcio da $10.4$ desintegraciones por minuto y por gramo. Dátalo.
        ¿Es verosímil un escriba de hace veintidós siglos?
  \item La \omterm{def:g11:nucleus-radioactivity:activity}{actividad} se mide con un margen de $\pm 0.2$ desintegraciones por minuto.
        Estima la incertidumbre de la edad del papiro usando
        $\Delta t \approx \Delta\mathcal A/(\lambda \mathcal A)$.
  \item El pergamino ``antiguo'' de un anticuario da $13.4$ desintegraciones por
        minuto y por gramo. ¿Veredicto?
\end{enumerate}

\medskip
\noindent\textbf{Parte III --- Jueves: los límites del reloj.}
\begin{enumerate}[resume]
  \item ¿Tras cuántos \omterm{def:g11:nucleus-radioactivity:halflife}{periodos de semidesintegración} cae la \omterm{def:g11:nucleus-radioactivity:activity}{actividad} de una muestra
        por debajo de las $0.2$ desintegraciones por minuto y por gramo de los
        contadores? ¿A qué edad corresponde eso?
  \item Los mejores laboratorios llegan a unos diez \omterm{def:g12:oscillators-and-time:oscillator}{periodos}. ¿Cuál es el horizonte
        práctico de la datación por radiocarbono, en años?
  \item Un hueso de dinosaurio tiene unos \num{6.6e7} años: ¿cuántos \omterm{def:g12:oscillators-and-time:oscillator}{periodos} del
        carbono-14 son? Usando la pregunta 4, ¿tras cuántos \omterm{def:g12:oscillators-and-time:oscillator}{periodos} no queda ni
        \emph{un} \omterm{def:g11:fundamental-interactions:ladder}{átomo} del carbono-14 de ese gramo? Concluye sobre eso de ``datar
        dinosaurios con carbono''.
  \item ¿Qué tipo de \omterm{def:g11:nucleus-radioactivity:notation}{nucleido} podría datar en cambio la roca del dinosaurio?
\end{enumerate}

\medskip
\noindent\textbf{Parte IV --- Viernes: la roca volcánica.}
\begin{enumerate}[resume]
  \item El argón es un gas y el potasio, un fiel residente de los sólidos: explica por
        qué la solidificación de la lava pone a cero el reloj potasio--argón.
  \item Suponiendo que cada \omterm{def:g11:fundamental-interactions:ladder}{núcleo} de potasio-40 desintegrado deja un \omterm{def:g11:fundamental-interactions:ladder}{átomo} de argón
        atrapado, demuestra que
        $N_{\mathrm{Ar}}/N_{\mathrm K} = e^{\lambda t} - 1$.
  \item La roca que hay debajo del fósil del museo presenta
        $N_{\mathrm{Ar}}/N_{\mathrm K} = 0.60$. Data la roca.
  \item ¿Podría el carbono-14 haber datado esta roca, o el potasio--argón la estatua?
        Da la ventana útil de cada reloj en \omterm{def:g11:nucleus-radioactivity:halflife}{periodos de semidesintegración}.
  \item Informe del viernes, en una frase: las tres edades halladas esta semana y la
        única ley que las leyó todas.
\end{enumerate}
\end{problem}
